\documentclass[12pt]{article}
\usepackage{amsmath,amssymb,graphicx}
\usepackage{geometry}
\geometry{margin=1in}

\title{Stablecoins, Parallel Exchange Rates, and Fragility in Fixed Exchange Rate Regimes}
\author{}
\date{}

\begin{document}
\maketitle

\begin{abstract}
We study how stablecoins affect parallel foreign exchange markets under fixed exchange rate regimes. Stablecoins act as both a transactional innovation and an endogenous public signal. We show that while stablecoins improve efficiency, they can increase fragility by facilitating coordination. Welfare is non-monotonic in exchange rate misalignment.
\end{abstract}

\input{sections/01_intro_part1}
\input{sections/01_intro_part2}
\input{sections/01_intro_part3}

\section{Model}

\subsection{Environment}

The economy consists of a continuum of agents indexed by $i \in [0,1]$. Time is static. There is an underlying fundamental $\theta \in \mathbb{R}$ representing exchange rate misalignment, where higher values of $\theta$ correspond to a greater overvaluation of the official exchange rate.

The fundamental $\theta$ is drawn from a normal prior distribution:
\begin{equation}
\theta \sim \mathcal{N}(\mu, \sigma_\theta^2).
\end{equation}

Agents are risk averse and derive utility from consumption. They face uncertainty about the true level of exchange rate misalignment and must form beliefs based on imperfect information.

\subsection{Motives for Foreign Currency Demand}

Agents demand foreign currency for two distinct reasons:

\begin{enumerate}
\item \textbf{Hedging motive:} Agents acquire foreign currency to insure against the risk of a devaluation.
\item \textbf{Speculative motive:} Agents acquire foreign currency to exit ahead of a potential regime collapse.
\end{enumerate}

These two motives generate different components of aggregate foreign currency demand, which play distinct roles in equilibrium.

\subsection{Assets}

Agents can access foreign currency through two instruments:

\paragraph{Cash foreign currency ($c$):}
\begin{itemize}
\item Involves higher transaction costs
\item Trades in decentralized and opaque markets
\item Provides only noisy, privately observed information
\end{itemize}

\paragraph{Stablecoins ($s$):}
\begin{itemize}
\item Lower transaction costs
\item Scalable and accessible through digital platforms
\item Generates a publicly observable market price
\end{itemize}

Let $q_c$ and $q_s$ denote holdings of cash and stablecoins, respectively. Total foreign currency holdings are given by:
\begin{equation}
q = q_c + q_s.
\end{equation}

\subsection{Transaction Costs}

The cost of acquiring foreign currency depends on the instrument used. We assume convex transaction costs:
\begin{align}
C_c(q_c) &= a_c + b_c \sqrt{q_c} + \kappa_c, \\
C_s(q_s) &= a_s + b_s \sqrt{q_s} + \kappa_s,
\end{align}

with $C_s(q_s) < C_c(q_c)$ for comparable quantities, reflecting the efficiency advantage of stablecoins.

\subsection{Discussion}

The two-asset structure is essential. It allows us to separate the transaction cost channel from the information channel. A one-asset model with improved information would confound these mechanisms. In contrast, the presence of both cash and stablecoins enables us to study how financial innovation affects both market efficiency and the information environment.

\subsection{Information Structure}

Agents do not observe the fundamental $\theta$ directly. Instead, they receive information from two sources corresponding to the two foreign currency instruments.

\paragraph{Private signal (cash price)}
Each agent $i$ observes a private signal:
\begin{equation}
x_i = \theta + \epsilon_i, \quad \epsilon_i \sim \mathcal{N}(0, \sigma_x^2),
\end{equation}
which can be interpreted as the price in decentralized cash markets. This signal is noisy and privately observed.

\paragraph{Public signal (stablecoin price)}
All agents observe a common public signal:
\begin{equation}
s = \theta + \eta, \quad \eta \sim \mathcal{N}(0, \sigma_s^2(q_s)).
\end{equation}

Crucially, the variance of the public signal depends on equilibrium stablecoin usage:
\begin{equation}
\sigma_s^2(q_s) = \frac{\bar{\sigma}_s^2}{1 + \rho q_s}, \quad \rho > 0.
\end{equation}

Thus, higher stablecoin adoption increases the precision of the public signal.

\subsection{Endogenous Information}

The dependence of $\sigma_s^2$ on $q_s$ introduces a feedback loop:
\begin{enumerate}
\item Higher $q_s$ \rightarrow more informative public signal
\item More informative signal \rightarrow stronger coordination
\item Stronger coordination \rightarrow changes in asset demand
\end{enumerate}

This endogeneity is central to the model and distinguishes it from standard global-games frameworks with exogenous public signals.

\subsection{Beliefs}

Agents form posterior beliefs about $\theta$ using Bayes' rule. Given normal priors and signals, the posterior mean is linear:
\begin{equation}
\mathbb{E}[\theta \mid x_i, s] = w_0 \mu + w_x x_i + w_s s,
\end{equation}
where weights depend on signal precisions:
\begin{align}
\tau_\theta &= \frac{1}{\sigma_\theta^2}, \\
\tau_x &= \frac{1}{\sigma_x^2}, \\
\tau_s(q_s) &= \frac{1}{\sigma_s^2(q_s)}.
\end{align}

The weights are given by:
\begin{align}
w_0 &= \frac{\tau_\theta}{\tau_\theta + \tau_x + \tau_s}, \\
w_x &= \frac{\tau_x}{\tau_\theta + \tau_x + \tau_s}, \\
w_s &= \frac{\tau_s}{\tau_\theta + \tau_x + \tau_s}.
\end{align}

\subsection{Implications}

As stablecoin usage increases, $\tau_s$ increases and agents place more weight on the public signal $s$. This shifts beliefs toward common information and reduces the dispersion of posterior beliefs across agents. As a result, agents' actions become more correlated, strengthening coordination in equilibrium.

\subsection{Decisions and Strategies}

Agents decide whether to engage in a speculative attack (i.e., acquire foreign currency for exit) based on their posterior beliefs about the fundamental $\theta$.

Following the global-games literature, we restrict attention to monotone strategies. There exists a cutoff $x^*$ such that agent $i$ attacks if and only if:
\begin{equation}
x_i \geq x^*.
\end{equation}

\subsection{Aggregate Attack}

Given the cutoff strategy, the measure of agents who attack is:
\begin{equation}
R(\theta) = \mathbb{P}(x_i \geq x^* \mid \theta) = 1 - \Phi\left(\frac{x^* - \theta}{\sigma_x}\right),
\end{equation}
where $\Phi(\cdot)$ is the standard normal cumulative distribution function.

\subsection{Hedging Demand}

In addition to speculative demand, agents demand foreign currency for hedging purposes. We model aggregate hedging demand as an increasing function of expected misalignment:
\begin{equation}
H = h \cdot \mathbb{E}[\theta \mid s], \quad h > 0.
\end{equation}

This captures the idea that even in the absence of coordination, agents acquire foreign currency as insurance against depreciation risk.

\subsection{Crisis Condition}

A crisis occurs if total demand for foreign currency exceeds the available policy buffer. We model this as:
\begin{equation}
R(\theta) + H \geq \alpha - \delta \theta,
\end{equation}
where:
\begin{itemize}
\item $\alpha$ represents the strength of policy buffers (e.g., reserves)
\item $\delta$ captures the fact that higher misalignment weakens the regime
\end{itemize}

\subsection{Crisis Probability}

Given the information structure, agents form beliefs about the probability of a crisis:
\begin{equation}
\pi(x_i, s) = \mathbb{P}\left(R(\theta) + H \geq \alpha - \delta \theta \mid x_i, s\right).
\end{equation}

\subsection{Payoffs}

Agents derive utility from consumption, which depends on whether a crisis occurs.

\paragraph{No-crisis consumption}
\begin{equation}
c^N = y + \zeta \log(1 + q_c + q_s) - C_c(q_c) - C_s(q_s),
\end{equation}

\paragraph{Crisis consumption}
\begin{equation}
c^C = c^N - \lambda \theta (1 + \xi \omega_s),
\end{equation}

where $\omega_s = \frac{q_s}{q_c + q_s}$ is the share of stablecoins in total foreign currency holdings.

\paragraph{Expected utility}
\begin{equation}
U = (1 - \pi) \log(c^N) + \pi \log(c^C).
\end{equation}

\subsection{Equilibrium Definition}

An equilibrium consists of:
\begin{itemize}
\item A cutoff $x^*$
\item Asset allocations $(q_c, q_s)$
\item Beliefs consistent with Bayesian updating
\item A crisis probability consistent with aggregate behavior
\end{itemize}

such that agents optimize given beliefs and the information structure, and all objects are mutually consistent.

\subsection{Key Mechanism}

Stablecoins affect equilibrium through two channels:
\begin{enumerate}
\item \textbf{Cost channel:} lower transaction costs increase foreign currency demand
\item \textbf{Information channel:} improved public signal strengthens coordination
\end{enumerate}

The second channel amplifies strategic complementarities and increases the likelihood of coordinated attacks, generating fragility in equilibrium.


\input{sections/03_theory_part1}
\subsection{Refined Equilibrium Characterization}

We make the cutoff structure explicit in terms of posterior beliefs. Let the posterior mean be:
\begin{equation}
\hat{\theta}_i = w_0 \mu + w_x x_i + w_s s.
\end{equation}

Since $\hat{\theta}_i$ is monotone in $x_i$, the cutoff strategy in signal space maps into a cutoff in posterior space. There exists $\hat{\theta}^*$ such that agents attack iff:
\begin{equation}
\hat{\theta}_i \geq \hat{\theta}^*.
\end{equation}

This implies a corresponding cutoff in signals:
\begin{equation}
x^* = \frac{\hat{\theta}^* - w_0 \mu - w_s s}{w_x}.
\end{equation}

\subsection{Sharpened Comparative Statics}

\textbf{Proposition 5 (Refined).} An increase in public signal precision $\tau_s$ lowers the posterior cutoff $\hat{\theta}^*$.

\textit{Proof.}

The indifference condition is:
\begin{equation}
\pi(\hat{\theta}^*) \cdot B = C.
\end{equation}

Higher $\tau_s$ increases the weight on the common signal and reduces posterior variance. This steepens the mapping from fundamentals to aggregate behavior. As a result, for a given benefit-cost ratio, the required posterior threshold for coordination decreases. Formally, $\frac{\partial \hat{\theta}^*}{\partial \tau_s} < 0$. \hfill $\square$

\textbf{Proposition 6 (Refined).} Public signal precision increases the sensitivity of aggregate attacks to fundamentals.

\textit{Proof.}

Aggregate attack is:
\begin{equation}
R(\theta) = 1 - \Phi\left(\frac{x^* - \theta}{\sigma_x}\right).
\end{equation}

Using the chain rule:
\begin{equation}
\frac{dR}{d\theta} = \phi\left(\frac{x^* - \theta}{\sigma_x}\right) \cdot \frac{1}{\sigma_x}.
\end{equation}

Since $x^*$ decreases in $\tau_s$, more mass lies near the threshold region, increasing local sensitivity. \hfill $\square$

\subsection{Global Games Mapping}

The model maps to a standard global games environment with endogenous public precision. The key sufficient statistic is total precision:
\begin{equation}
\tau = \tau_x + \tau_s(q_s).
\end{equation}

As $\tau_s$ increases, the dispersion of posterior beliefs shrinks:
\begin{equation}
\text{Var}(\hat{\theta}_i) = \frac{1}{\tau_x + \tau_s + \tau_\theta}.
\end{equation}

Lower dispersion implies tighter clustering of actions, strengthening coordination.

\subsection{Welfare Threshold (Refined)}

\textbf{Proposition 8 (Refined).} There exists a unique threshold $\theta^*$ such that:
\begin{equation}
\frac{dW}{dq_s} > 0 \quad \text{for } \theta < \theta^*, \quad \frac{dW}{dq_s} < 0 \text{ for } \theta > \theta^*.
\end{equation}

\textit{Proof.}

Differentiate welfare:
\begin{equation}
\frac{dW}{dq_s} = \frac{\partial W^{eff}}{\partial q_s} - \frac{\partial W^{frag}}{\partial q_s}.
\end{equation}

The efficiency term is decreasing in transaction costs and always positive. The fragility term is increasing in crisis probability, which is increasing in $\tau_s(q_s)$ and thus $q_s$. At low $\theta$, $\frac{\partial W^{frag}}{\partial q_s}$ is small; at high $\theta$, it dominates. Continuity ensures a unique crossing. \hfill $\square$

\subsection{Interpretation}

The refined results make clear that the core mechanism operates through the compression of beliefs. Stablecoins reduce disagreement, which is beneficial for price discovery but destabilizing in environments with strategic complementarities.


\input{sections/04_quant_part1}
\subsection{Quantitative Results}

This section presents the main quantitative results of the model and connects them to the theoretical predictions.

\subsubsection{Result 1: Stablecoin Demand and Misalignment}

Figure 1 shows stablecoin demand $q_s$ as a function of exchange rate misalignment $\theta$. Stablecoin usage increases monotonically with $\theta$.

This reflects the increasing demand for foreign currency as misalignment grows. Since stablecoins offer lower transaction costs, they capture a growing share of foreign currency demand.

\subsubsection{Result 2: Crisis Probability}

Figure 2 plots the probability of crisis as a function of $\theta$ under two scenarios: (i) cash-only economy and (ii) economy with stablecoins.

The presence of stablecoins increases crisis probability for all levels of misalignment. The difference is modest at low $\theta$ but becomes more pronounced as misalignment increases.

This result reflects the information channel: stablecoins improve signal precision, facilitating coordination and increasing the likelihood of speculative attacks.

\subsubsection{Result 3: Welfare}

Figure 3 shows welfare as a function of $\theta$. Welfare exhibits a non-monotonic relationship with misalignment.

At low levels of $\theta$, welfare is higher in the presence of stablecoins due to lower transaction costs and improved access to foreign currency. At higher levels of $\theta$, welfare is lower due to increased crisis probability.

\subsubsection{Result 4: Welfare Difference}

Figure 4 plots the difference in welfare between the stablecoin and cash-only economies.

The welfare difference is positive at low misalignment and negative at high misalignment, crossing zero at an intermediate threshold. This confirms the theoretical prediction of a welfare tradeoff.

\subsubsection{Magnitude and Realism}

Importantly, the magnitudes of the effects are moderate. Crisis probabilities remain well below one, and welfare changes are economically meaningful but not extreme.

This ensures that the results are consistent with realistic macro-financial environments rather than driven by extreme parameter choices.

\subsubsection{Link to Theory}

The quantitative results closely match the theoretical predictions:

\begin{itemize}
\item Increasing $q_s$ raises signal precision (Proposition 2)
\item Higher precision increases coordination (Propositions 5--7)
\item Welfare reflects the tradeoff between efficiency and fragility (Proposition 8)
\end{itemize}

Thus, the simulations provide a quantitative validation of the theoretical mechanism.

\input{sections/04_quant_part3}

\section{Figures}

\begin{figure}[h!]
\centering
\includegraphics[width=0.8\textwidth]{figures/fig1_asset_allocation.png}
\caption{Stablecoin demand increases with exchange rate misalignment}
\end{figure}

\begin{figure}[h!]
\centering
\includegraphics[width=0.8\textwidth]{figures/fig2_crisis.png}
\caption{Crisis probability with and without stablecoins}
\end{figure}

\begin{figure}[h!]
\centering
\includegraphics[width=0.8\textwidth]{figures/fig3_welfare.png}
\caption{Welfare is non-monotonic in misalignment}
\end{figure}

\section{Conclusion}

Stablecoins transform parallel exchange rate markets by improving both transaction efficiency and information aggregation. While these innovations generate welfare gains at low levels of misalignment, they can increase fragility by facilitating coordination at higher levels of misalignment. These results highlight a fundamental tradeoff between efficiency and stability in financial innovation.

\end{document}
